%%% FIELD def-trace-invariant-proposed
For a raw state \(\kappa=(\mathcal H,F,c,q,\pi,\Lambda)\), the predicate \(\operatorname{Inv}(\kappa)\) holds when all of the following conditions hold on one common support. First, \(q\) is a normalized kernel that is nested Markov with respect to the current CADMG \(\mathcal H\). Second, every denominator already used in the operation word and every denominator of a legal continuation is nonzero on that support. Third, for every compatible specialization of the observed-cell coordinates, applying any such supported continuation to the rational provenance expression \(\pi\) gives the same kernel as applying the identical ordered nested-kernel operations to \(q\). Fourth, each entry \((\lambda_i,\pi_i)\) of \(\Lambda\) is the exact pre-operation provenance expression at which \(\lambda_i\) was valid. This definition names the invariant; it does not assert that delayed context selection preserves it.

%%% FIELD def-trace-state-current
\(\kappa=(\mathcal H,F,c,q,\pi,\Lambda)\), where \(\mathcal H\) is the current CADMG, \(F\) is its fixed set, \(c\) is a partial assignment to missingness indicators, \(q\) is a normalized kernel on the random vertices of \(\mathcal H\), \(\pi\) is a finite expression in observed cells of \(P_{\mathrm{obs}}\), and \(\Lambda\) is an ordered list of pairs consisting of a label license and the pre-restriction trace that licensed it.

%%% FIELD def-trace-state-proposed
A raw trace-aware colored state is a tuple \(\kappa=(\mathcal H,F,c,q,\pi,\Lambda)\), where \(\mathcal H\) is the current CADMG, \(F\) is its fixed set, \(c\) is a partial assignment to missingness indicators, \(q\) is a normalized kernel on the random vertices of \(\mathcal H\), \(\pi\) is a finite rational expression in observed cells of \(P_{\mathrm{obs}}\), and \(\Lambda\) is an ordered list of pairs consisting of a label license and its exact pre-operation provenance expression. A raw state is called certified when \(\operatorname{Inv}(\kappa)\) holds. Normalization of \(q\) alone does not certify Markovness or provenance realization.

%%% FIELD def-colored-fixing-current
If \(v\) is fixable in \(\mathcal H\) and license \(\lambda\) is valid on the trace recorded by \(\kappa\), \(\Phi_v^{\lambda}\) moves \(v\) from the random to the fixed set, deletes arrowheads into \(v\), replaces \(q\) by \(q/q(v\mid\operatorname{mb}_{\mathcal H}(v))\), applies the same division to \(\pi\), and appends \(\lambda\) and its current provenance to \(\Lambda\).

%%% FIELD def-colored-fixing-proposed
Let \(d_v=q(v\mid\operatorname{mb}_{\mathcal H}(v))\). If \(v\) is fixable in \(\mathcal H\), license \(\lambda\) is valid on the current trace, and \(d_v\) is nonzero on the stated support, the raw update \(\Phi_v^{\lambda}\) moves \(v\) from the random to the fixed set, deletes arrowheads into \(v\), and replaces \(q\) by \(q/d_v\). When \(\operatorname{Inv}(\kappa)\) supplies a rational provenance denominator \(d_v^{\pi}\) whose compatible specializations equal \(d_v\), the update replaces \(\pi\) by \(\pi/d_v^{\pi}\) and appends \((\lambda,\pi)\) to \(\Lambda\). The update definition does not by itself assert that its output satisfies \(\operatorname{Inv}\).

%%% FIELD def-context-selection-current
For a missingness assignment \(c'\) compatible with the current trace, \(\Sigma_{c'}(\kappa)\) restricts \(q\) and \(\pi\) to \(c'\), normalizes on the positive selected cell, deletes the edges carrying labels activated by \(c'\), and retains every earlier entry of \(\Lambda\).

%%% FIELD def-context-selection-proposed
For a missingness assignment \(c'\) compatible with the current trace and having positive selected probability, the raw update \(\Sigma_{c'}(\kappa)\) conditions \(q\) and \(\pi\) on \(c'\), divides each by its displayed positive normalizing constant, deletes the edges carrying labels activated by \(c'\), and retains every earlier entry of \(\Lambda\). This syntactic update does not assert that the conditioned kernel is Markov with respect to the edge-deleted CADMG or that \(\operatorname{Inv}\) is preserved; those are separate proof obligations.

%%% FIELD def-adjacent-swap-current
\(\operatorname{Swap}_{\kappa}(u,v)\) holds exactly when \(u\) and \(v\) are both fixable in the current \(\mathcal H\), \(u\notin\operatorname{an}_{\mathcal H}(C_v)\), \(v\notin\operatorname{an}_{\mathcal H}(C_u)\), the license contexts \(C_u\) and \(C_v\) agree on their common indicators, and both fixing denominators are supported in the same pre-restriction trace.

%%% FIELD def-adjacent-swap-proposed
\(\operatorname{Swap}_{\kappa}(u,v)\) holds exactly when \(\operatorname{Inv}(\kappa)\) holds, \(u\) and \(v\) are both fixable in the current \(\mathcal H\), \(u\notin\operatorname{an}_{\mathcal H}(C_v)\), \(v\notin\operatorname{an}_{\mathcal H}(C_u)\), the license contexts \(C_u\) and \(C_v\) agree on their common indicators, and both fixing denominators are supported in the same pre-restriction trace. This predicate is only a conditional adjacent-swap test; it does not assert a global preservation or confluence theorem.

%%% FIELD def-figure-current
\(\mathcal G_{\mathrm{3b}}\) has substantive vertices \(\{A,Y,Z,W\}\), indicators \(\{R_W,R_Z\}\), their observed proxies, arrows \(R_W\to A\), \(Z\to A\), \(Z\to W\), \(W\to A\), \(A\to Y\), edges \(W\leftrightarrow Y\) and \(R_Z\leftrightarrow Y\), and label \(R_W=0\) on \(W\to A\); \(P_{\mathrm{3b}}\) is any compatible strictly positive observed law.

%%% FIELD def-figure-proposed
\(\mathcal G_{\mathrm{3b}}\) has substantive vertices \(\{A,Y,Z,W\}\), indicators \(\{R_W,R_Z\}\), their observed proxies, arrows \(R_W\to A\), \(Z\to A\), \(Z\to W\), \(W\to A\), \(A\to Y\), edges \(W\leftrightarrow Y\) and \(R_Z\leftrightarrow Y\), and label \(R_W=0\) on \(W\to A\); \(P_{\mathrm{3b}}\) is any compatible strictly positive observed law. This graph-and-law definition does not itself designate a legal colored operation word, a fixing denominator, or a pre-selection provenance certificate.

%%% FIELD def-colored-hedge-current
\(\operatorname{CH}(\mathcal G,A,Y)\) is the set of tuples \((\mathcal F,\mathcal F',\chi,T)\) such that \(\mathcal F'\subseteq\mathcal F\) are equally rooted \(C\)-forests, \(\mathcal F\cap A\neq\varnothing\), \(\mathcal F'\cap A=\varnothing\), their roots are ancestors of \(Y\) after intervention on \(A\), \(\chi\) assigns every label-sensitive edge its missingness context, and \(T\) is a prefix-closed collection of observed traces on which every color used by the two forests is licensed.

%%% FIELD def-colored-hedge-proposed
\(\operatorname{CH}(\mathcal G,A,Y)\) is the finite combinatorial candidate set of tuples \((\mathcal F,\mathcal F',\chi,T)\) such that \(\mathcal F'\subseteq\mathcal F\) are equally rooted \(C\)-forests, \(\mathcal F\cap A\neq\varnothing\), \(\mathcal F'\cap A=\varnothing\), their common roots are ancestors of \(Y\) after intervention on \(A\), \(\chi\) assigns a compatible missingness context to every label-sensitive forest edge, and \(T\) is a finite prefix-closed collection of observed traces on which every assigned color is licensed. Membership is purely graphical and trace-syntactic: it does not assert that a graph-only recursion returns the tuple or that two strictly positive lm-SCMs realize it.

%%% FIELD def-compiler-current
Given an exact specialized negative-cell witness \(\omega=(a,y,\theta^0,\theta^1)\) and a finite retained operation history \(\mathcal T\), first decode \(\theta^0,\theta^1\) through their selected response tables and product laws over the individual edge and private factors, obtaining \(\mathcal M_0,\mathcal M_1\). Next enumerate, in a fixed prefix and vertex-subset order, every triple \((\kappa_*,D,D')\) occurring in \(\mathcal T\) and test the conjuncts of the candidate predicate \(\Delta(\omega,\mathcal T)=(\kappa_*,D,D')\). Trace equalities are tested after clearing the supported nonzero denominators, using the Thom encodings supplied by exact specialization. Return \((\operatorname{found}(\kappa_*,D,D'),\mathcal M_0,\mathcal M_1)\) for the first successful test and return \((\operatorname{exhausted},\mathcal M_0,\mathcal M_1)\) if the finite enumeration ends. This is a partial graphical diagnostic: neither tag asserts existence of a colored hedge, membership in \(\operatorname{CH}(\mathcal G,A,Y)\), realization of a forest by the decoded models, or alignment of the algebraic sample with a graphwise recursion failure.

%%% FIELD def-compiler-proposed
Given an exact specialized negative coordinate-cell witness \(\omega=(a,y,\theta^0,\theta^1)\) and a finite retained raw operation history \(\mathcal T\), first decode \(\theta^0,\theta^1\) through their selected response tables and product laws over the individual edge and private factors, obtaining \(\mathcal M_0,\mathcal M_1\). Next enumerate, in a fixed prefix and vertex-subset order, every triple \((\kappa_*,D,D')\) occurring in \(\mathcal T\) and test the conjuncts of the candidate predicate \(\Delta(\omega,\mathcal T)=(\kappa_*,D,D')\). Trace equalities are tested after clearing the displayed supported nonzero denominators, using the Thom encodings supplied by exact specialization. Return \((\operatorname{found}(\kappa_*,D,D'),\mathcal M_0,\mathcal M_1)\) for the first successful test and return \((\operatorname{exhausted},\mathcal M_0,\mathcal M_1)\) if the finite enumeration ends. This diagnostic is independent of any graph-only recursion. Neither tag asserts \(\operatorname{Inv}(\kappa_*)\), existence or membership of a colored hedge, realization of a forest by the decoded models, or alignment of the algebraic sample with a graphwise recursion failure.

%%% FIELD def-tcf-current
On formal input \((\mathcal G,A,Y,\mathbf p)\), first construct the response-polynomial feasibility formula \(\exists\theta\,\mathsf C_{\mathcal G}(\theta;\mathbf p)\), all coordinate-constancy formulas \(\mathsf I_{a,y}(\mathcal G;\mathbf p)\), and the direct target-constancy formula \(\mathsf I_{\tau}(\mathcal G;\mathbf p,\mathbf s)\) with \(\mathbf s\) retained as a formal real vector. Run one signed-subresultant CAD for the coordinate family in \(\mathbf p\), and a second signed-subresultant CAD for the target family in \((\mathbf p,\mathbf s)\) together with the lifted coordinate-cell polynomials. This produces the canonical coordinate cells \(\psi_j\) and target cells \(\chi_k\), with their respective unique-value or witness relations, before any fixing trace is enumerated.

After both base partitions have been fixed, enumerate expression-discovery traces. Start from every supported observed margin and partial-context trace kernel written as a rational expression in \(\mathbf p\). Recursively branch only by strict ancestral marginalization, restriction to a proper child of a nontrivial district decomposition, a licensed \(\Phi_v^{\lambda}\) update, or a delayed \(\Sigma_{c'}\) update that strictly extends the partial context. Retain the complete operation word and memoize only literal complete tuples \(\kappa=(\mathcal H,F,c,q,\pi,\Lambda)\), including the syntax of \(q\), \(\pi\), and every provenance entry in \(\Lambda\). An adjacent-swap optimization may merge branches only after replaying both license orders from their common pre-swap state in a fixed canonical order and verifying literal equality of the complete resulting tuples. Fixing and selection are never commuted.

For a candidate coordinate expression \(r(\mathbf p)=N(\mathbf p)/D(\mathbf p)\), a coordinate cell \(\psi_j\), and \((a,y)\) having a unique relation, decide
\[
\forall\mathbf p\,\forall z
\left[
\psi_j(\mathbf p)\wedge\mathsf V_{a,y}(\mathbf p,z)
\Longrightarrow
D(\mathbf p)\neq0\wedge D(\mathbf p)z=N(\mathbf p)
\right].
\]
For a candidate direct-target expression \(r_{\tau}(\mathbf p,\mathbf s)=N_{\tau}(\mathbf p,\mathbf s)/D_{\tau}(\mathbf p,\mathbf s)\), formed as a score-linear combination of compatible trace expressions, and a target cell \(\chi_k\) having a unique relation, decide
\[
\forall\mathbf p\,\forall\mathbf s\,\forall z
\left[
\chi_k(\mathbf p,\mathbf s)\wedge\mathsf V_{\tau}(\mathbf p,\mathbf s,z)
\Longrightarrow
D_{\tau}(\mathbf p,\mathbf s)\neq0
\wedge D_{\tau}(\mathbf p,\mathbf s)z=N_{\tau}(\mathbf p,\mathbf s)
\right].
\]
Annotate a cell only when its implication is true. If no candidate passes, retain the corresponding unique-value graph; on a nonconstant coordinate or target cell retain \(\mathsf W_j\) or \(\mathsf W_{\tau}\), respectively. Trace enumeration and annotation never split, merge, or otherwise alter either base partition. The notation \(\operatorname{TCF}(\mathcal G,A,Y,P_{\mathrm{obs}})\) denotes the two symbolic partitions and their verified annotations, with formal names \(p_o=P_{\mathrm{obs}}(o)\); a supplied \(s_Y\) selects the target-score section only at mathematical or exact specialization. The procedure neither names nor returns a colored hedge.

%%% FIELD def-tcf-proposed
On formal input \((\mathcal G,A,Y,\mathbf p)\), first construct the response-polynomial feasibility formula \(\exists\theta\,\mathsf C_{\mathcal G}(\theta;\mathbf p)\), all coordinate-constancy formulas \(\mathsf I_{a,y}(\mathcal G;\mathbf p)\), and the direct target-constancy formula \(\mathsf I_{\tau}(\mathcal G;\mathbf p,\mathbf s)\) with \(\mathbf s\) retained as a formal real vector. Run one signed-subresultant CAD for the coordinate family in \(\mathbf p\), and a second signed-subresultant CAD for the target family in \((\mathbf p,\mathbf s)\) together with the lifted coordinate-cell polynomials. This produces the canonical coordinate cells \(\psi_j\) and target cells \(\chi_k\), with their respective unique-value or witness relations, before any fixing trace is enumerated.

After both base partitions have been fixed, enumerate raw expression-discovery traces. Start from every supported observed margin and partial-context trace kernel written as a rational expression in \(\mathbf p\). Recursively branch only by strict ancestral marginalization, restriction to a proper child of a nontrivial district decomposition, a licensed \(\Phi_v^{\lambda}\) update, or a delayed \(\Sigma_{c'}\) update that strictly extends the partial context. Retain the complete operation word and memoize only literal complete tuples \(\kappa=(\mathcal H,F,c,q,\pi,\Lambda)\), including the syntax of \(q\), \(\pi\), and every provenance entry in \(\Lambda\). Distinct operation words are never merged, even when an adjacent-swap audit proves equality after specialization; canonical replay may be recorded as an audit annotation only. Fixing and selection are never commuted. The enumeration does not assume that every raw state satisfies \(\operatorname{Inv}\).

For a candidate coordinate expression \(r(\mathbf p)=N(\mathbf p)/D(\mathbf p)\), a coordinate cell \(\psi_j\), and \((a,y)\) having a unique relation, decide
\[
\forall\mathbf p\,\forall z
\left[
\psi_j(\mathbf p)\wedge\mathsf V_{a,y}(\mathbf p,z)
\Longrightarrow
D(\mathbf p)\neq0\wedge D(\mathbf p)z=N(\mathbf p)
\right].
\]
For a candidate direct-target expression \(r_{\tau}(\mathbf p,\mathbf s)=N_{\tau}(\mathbf p,\mathbf s)/D_{\tau}(\mathbf p,\mathbf s)\), formed as a score-linear combination of compatible trace expressions, and a target cell \(\chi_k\) having a unique relation, decide
\[
\forall\mathbf p\,\forall\mathbf s\,\forall z
\left[
\chi_k(\mathbf p,\mathbf s)\wedge\mathsf V_{\tau}(\mathbf p,\mathbf s,z)
\Longrightarrow
D_{\tau}(\mathbf p,\mathbf s)\neq0
\wedge D_{\tau}(\mathbf p,\mathbf s)z=N_{\tau}(\mathbf p,\mathbf s)
\right].
\]
Annotate a cell only when its implication is true. This whole-cell implication, rather than an unproved trace invariant, is the soundness gate for every rational annotation. If no candidate passes, retain the corresponding unique-value graph; on a nonconstant coordinate or target cell retain \(\mathsf W_j\) or \(\mathsf W_{\tau}\), respectively. Trace enumeration and annotation never split, merge, or otherwise alter either base partition. The notation \(\operatorname{TCF}(\mathcal G,A,Y,P_{\mathrm{obs}})\) denotes the two symbolic partitions and their verified annotations, with formal names \(p_o=P_{\mathrm{obs}}(o)\); a supplied \(s_Y\) selects the target-score section only at mathematical or exact specialization. The procedure neither names nor returns a colored hedge and is not a graph-only sound-and-complete recursion.

%%% FIELD swap-current-statement
Suppose \(\operatorname{Swap}_{\kappa}(u,v)\) holds and both composite colored fixings are legal. After erasing the ordered provenance list \(\Lambda\), the two resulting states have the same CADMG, fixed set, context, normalized kernel, and observed-law provenance. Their two appended trace suffixes are respectively \(\bigl((\lambda_u,\pi),(\lambda_v,\Phi_u^{\lambda_u}\pi)\bigr)\) and \(\bigl((\lambda_v,\pi),(\lambda_u,\Phi_v^{\lambda_v}\pi)\bigr)\); hence the definition does not imply equality of the full trace states up to permutation. Branches may be merged only after replaying the two licenses in a fixed canonical order from their common pre-swap state.

%%% FIELD swap-proposed-statement
Suppose \(\operatorname{Inv}(\kappa)\), \(\operatorname{Swap}_{\kappa}(u,v)\), and legality of both composite colored fixings. The two orders have the same CADMG, fixed set, context, and normalized kernel. At every compatible observed-law specialization on their common support, their final provenance expressions evaluate to that same kernel. Their appended suffixes are respectively \(\bigl((\lambda_u,\pi),(\lambda_v,\Phi_u^{\lambda_u}\pi)\bigr)\) and \(\bigl((\lambda_v,\pi),(\lambda_u,\Phi_v^{\lambda_v}\pi)\bigr)\). Thus neither literal equality of the complete trace states nor global legal-order confluence follows; \(\operatorname{TCF}\) retains both branches and uses canonical replay only as an audit.

%%% FIELD swap-proof
Write the common initial state as
\[
\kappa=(\mathcal H,F,c,q,\pi,\Lambda).
\tag{1}
\]
By \(\operatorname{Inv}(\kappa)\), the kernel \(q\) is Markov with respect to \(\mathcal H\) on the common support. By \(\operatorname{Swap}_{\kappa}(u,v)\), both vertices are fixable in \(\mathcal H\), and the two fixing denominators are nonzero on that support. Legality of both composites supplies the corresponding intermediate domains. These are precisely the kernel hypotheses needed by \(\text{lem:ordinary-adjacent-fixing-commutation}\); hence
\[
\phi_v\{\phi_u(q;\mathcal H);\phi_u(\mathcal H)\}
=
\phi_u\{\phi_v(q;\mathcal H);\phi_v(\mathcal H)\},
\tag{2}
\]
and the two ordinary orders have the same final CADMG. Each order adds exactly \(u\) and \(v\) to \(F\), while neither colored fixing changes \(c\). This proves equality of the CADMG, fixed-set, context, and normalized-kernel components.

Fix any compatible specialization of the observed cells at which the recorded denominators are nonzero. The continuation clause in \(\operatorname{Inv}(\kappa)\) says that applying either ordered pair of supported divisions to \(\pi\) evaluates to applying the identical ordered pair to \(q\). Equation (2) therefore gives
\[
\left(\Phi_v^{\lambda_v}\Phi_u^{\lambda_u}\pi\right)(P_{\mathrm{obs}})
=
\left(\Phi_u^{\lambda_u}\Phi_v^{\lambda_v}\pi\right)(P_{\mathrm{obs}}),
\tag{3}
\]
and both sides specialize to the common kernel in (2). No equality of the two rational-expression strings is used.

Finally, the append rule in \(\text{def:colored-fixing}\) gives, in the order \(u,v\),
\[
\bigl((\lambda_u,\pi),(\lambda_v,\Phi_u^{\lambda_u}\pi)\bigr),
\tag{4}
\]
whereas the order \(v,u\) gives
\[
\bigl((\lambda_v,\pi),(\lambda_u,\Phi_v^{\lambda_v}\pi)\bigr).
\tag{5}
\]
The second provenance expressions in (4) and (5) are generally different strings, and permuting the list positions does not identify them. Thus (2)--(3) imply neither literal equality of complete trace states nor confluence of arbitrary legal words. The corrected \(\text{def:tcf-procedure}\) retains both words and allows canonical replay only as an audit, which proves the final clause.

%%% FIELD figure-current-statement
For every compatible \(P_{\mathrm{3b}}\), a legal trace on \(\mathcal G_{\mathrm{3b}}\) first derives \(P(Y\mid\operatorname{do}(A=a),R_W)\) and then applies \(\Sigma_{R_W=0}\), yielding \[P(Y\mid\operatorname{do}(A=a))=P(Y\mid\operatorname{do}(A=a),R_W=0)=P_{\mathrm{3b}}(Y\mid A=a,R_W=0).\] Hence the unique TCF partition cell containing \(P_{\mathrm{3b}}\) carries this functional. The closure generated by restricting to \(R_W=0\) before any fixing step contains no kernel with the pre-selection provenance needed for the first equality.

%%% FIELD figure-proposed-statement
For every \(\mathcal M\in\mathfrak M_{+}^{\mathrm{rm}}(\mathcal G_{\mathrm{3b}})\) inducing \(P_{\mathrm{3b}}\), every \(a\in\mathcal X_A\), and every \(y\in\mathcal X_Y\),
\[
Q_a^{\mathcal M}(y)
=P_{\mathcal M}(Y=y\mid\operatorname{do}(A=a),R_W=0)
=P_{\mathrm{3b}}(Y=y\mid A=a,R_W=0).
\]
All displayed conditioning events have positive probability. This proposition establishes the Figure 3(b) observed-law identity only; it does not assert that the frozen colored rules generate a complete legal fixing word, that \(\operatorname{TCF}\) enumerates that word, or that the restrict-first closure lacks an identifying expression.

%%% FIELD figure-proof
Fix \(\mathcal M\in\mathfrak M_{+}^{\mathrm{rm}}(\mathcal G_{\mathrm{3b}})\) inducing \(P_{\mathrm{3b}}\), and define \(Y(a)\) by replacing the structural equation for \(A\) by the constant \(a\). By \(\text{def:figure3b-witness}\), the only directed path from \(R_W\) to \(Y\) is \(R_W\to A\to Y\), and \(R_W\) is incident to no bidirected edge. The intervention deletes the arrow into \(A\), so the private exogenous factor determining \(R_W\) is disjoint from every exogenous factor determining \(Y(a)\). Mutual independence in \(\text{ass:semimarkovian}\) therefore gives
\[
Y(a)\mathbin{\perp\!\!\!\perp}R_W.
\tag{1}
\]

In the context \(R_W=0\), \(\text{ass:label-semantics}\) removes the input \(W\) from the equation for \(A\). Thus, conditional on \(R_W=0\), the natural treatment is a measurable function only of \(Z\), the private factor of \(A\), and the fixed context value. The counterfactual \(Y(a)\) is a function of \(a\), the private factor of \(Y\), and the individual factors on \(W\leftrightarrow Y\) and \(R_Z\leftrightarrow Y\). None of those factors occurs in the context-specific equation for \(A\), and the factor for \(Z\) is private. Factorwise independence consequently yields
\[
Y(a)\mathbin{\perp\!\!\!\perp}A\mid R_W=0.
\tag{2}
\]
No division has yet been used.

Strict full-data positivity in \(\text{ass:strict-positivity}\), together with the finite alphabets, implies
\[
P_{\mathcal M}(R_W=0)>0,
\qquad
P_{\mathcal M}(A=a,R_W=0)>0,
\tag{3}
\]
because each probability is a finite sum of positive full-data cells. Hence every following conditional is defined. Recursive evaluation of the same acyclic structural equations gives consistency, \(Y=Y(a)\) on \(\{A=a\}\). Using consistency, (2), the structural definition of intervention, and then (1), respectively, gives
\[
\begin{aligned}
P_{\mathcal M}(Y=y\mid A=a,R_W=0)
&=P_{\mathcal M}(Y(a)=y\mid A=a,R_W=0)\\
&=P_{\mathcal M}(Y(a)=y\mid R_W=0)\\
&=P_{\mathcal M}(Y=y\mid\operatorname{do}(A=a),R_W=0)\\
&=P_{\mathcal M}(Y=y\mid\operatorname{do}(A=a))\\
&=Q_a^{\mathcal M}(y).
\end{aligned}
\tag{4}
\]
The variables \(A,Y,R_W\) are observed in the Figure 3(b) record, and \(\text{ass:proxy-consistency}\) identifies every proxy coordinate used to form that record. Since \(\mathcal M\) induces \(P_{\mathrm{3b}}\), the first term of (4) equals \(P_{\mathrm{3b}}(Y=y\mid A=a,R_W=0)\). This proves the displayed identity and the support assertion. No colored transition, denominator license, enumeration, or restrict-first closure claim was invoked, which proves the stated scope limitation.

%%% FIELD response-fiber-proof
Fix a compatible observed law \(P_{\mathrm{obs}}\). First let \(\mathcal M\in\mathfrak F(\mathcal G,P_{\mathrm{obs}})\). By \(\text{def:observational-fiber}\) and \(\text{def:positive-rm-class}\), the endogenous alphabets are finite, the individual-edge and private exogenous factors are mutually independent, and every full-data cell is positive. The already established \(\text{lem:finite-factor-support-reduction}\) applies to this model and gives a label-compatible deterministic response representation in which each factor has support at most
\[
K_{\mathrm{ub}}(\mathcal G)
=\left|\mathcal X_{\mathcal V}\times\mathcal X_{\mathcal R}\right|
+|\mathcal X_A|(|\mathcal X_Y|-1),
\tag{1}
\]
the factors remain independent, and simultaneously
\[
P(V=v,R=r)=P_{\mathcal M}(v,r),
\qquad
P(Y=y\mid\operatorname{do}(A=a))=Q_a^{\mathcal M}(y)
\tag{2}
\]
for every \((v,r,a,y)\). Pad a smaller factor support by zero-weight states. Let \(\rho_{h,k}\) be the resulting factor weights and let \(g_*\in\mathcal T_{\mathcal G}\) be the tuple of restricted deterministic tables. Label compatibility follows from \(\text{ass:label-semantics}\). Set
\[
\gamma_{g_*}=1,
\qquad
\gamma_g=0\quad(g\ne g_*).
\tag{3}
\]
Equations (3) and the factor simplices give
\[
\gamma_g^2-\gamma_g=0,
\quad \sum_g\gamma_g=1,
\quad \rho_{h,k}\geq0,
\quad \sum_k\rho_{h,k}=1.
\tag{4}
\]

By mutual independence in \(\text{ass:semimarkovian}\), deterministic evaluation of the response tables yields
\[
\begin{aligned}
p_{v,r}(\theta_{\mathcal M})
&=\sum_u\mathbf 1\{g_*(u)=(v,r)\}\prod_h\rho_{h,u_h}
=P_{\mathcal M}(v,r),\\
q_{a,y}(\theta_{\mathcal M})
&=Q_a^{\mathcal M}(y).
\end{aligned}
\tag{5}
\]
Proxy consistency then gives, for every observed cell \(o\),
\[
p_o(\theta_{\mathcal M})
=\sum_{(v,r):\operatorname{proxy}(v,r)=o}P_{\mathcal M}(v,r)
=P_{\mathrm{obs}}(o).
\tag{6}
\]
The first equality in (5) and strict positivity give \(p_{v,r}(\theta_{\mathcal M})>0\). Equations (4)--(6) are exactly \(\mathsf C_{\mathcal G}(\theta_{\mathcal M};P_{\mathrm{obs}})\), proving the forward direction.

Conversely, suppose \(\mathsf C_{\mathcal G}(\theta;P_{\mathrm{obs}})\). Over \(\mathbb R\), \(\gamma_g^2-\gamma_g=0\) implies \(\gamma_g\in\{0,1\}\); together with \(\sum_g\gamma_g=1\), this selects a unique \(g_*\). For each factor index \(h\), define an exogenous variable \(U_h\) by
\[
P(U_h=k)=\rho_{h,k},
\qquad 1\leq k\leq K_{\mathrm{ub}}(\mathcal G),
\tag{7}
\]
and take the finitely many \(U_h\)'s mutually independent. Use the selected tuple \(g_*\) as the structural equations. Because \(g_*\in\mathcal T_{\mathcal G}\), these equations use only graph parents and satisfy every table equality imposed by a label. Together with the prescribed proxy equations, \(\text{def:published-lm-class}\), \(\text{ass:semimarkovian}\), \(\text{ass:label-semantics}\), and \(\text{ass:proxy-consistency}\) therefore give a compatible lm-SCM \(\mathcal M_\theta\).

Its product expansion is
\[
P_{\mathcal M_\theta}(v,r)
=\sum_u\mathbf 1\{g_*(u)=(v,r)\}\prod_h\rho_{h,u_h}
=p_{v,r}(\theta)>0.
\tag{8}
\]
The observed-cell conjuncts of \(\mathsf C_{\mathcal G}\) and proxy consistency imply
\[
P_{\mathcal M_\theta}(\mathcal O=o)
=p_o(\theta)=P_{\mathrm{obs}}(o).
\tag{9}
\]
Regularity and maximality are properties of the fixed input graph, finiteness is explicit in (7), and (8) is strict full-data positivity. Hence \(\mathcal M_\theta\in\mathfrak F(\mathcal G,P_{\mathrm{obs}})\). Finally, replacing the selected equation for \(A\) by the constant \(a\) in the same product expansion gives
\[
Q_a^{\mathcal M_\theta}(y)=q_{a,y}(\theta).
\tag{10}
\]
Equations (1)--(10) prove the two-way equality without replacing the individual-edge factorization by a district-level latent variable.

%%% FIELD tcf-soundness-proof
Let \(\mathbf p_{\mathrm{obs}}=(P_{\mathrm{obs}}(o):o\in\mathcal X_{\mathcal O})\), fix an identifying cell \(\psi_j\) containing \(\mathbf p_{\mathrm{obs}}\), and take \(\mathcal M\in\mathfrak F(\mathcal G,P_{\mathrm{obs}})\). By \(\text{thm:response-fiber-equivalence}\), there is \(\theta_{\mathcal M}\) such that
\[
\mathsf C_{\mathcal G}(\theta_{\mathcal M};\mathbf p_{\mathrm{obs}}),
\qquad
q_{a,y}(\theta_{\mathcal M})=Q_a^{\mathcal M}(y)
\tag{1}
\]
for every \((a,y)\). Fix one such pair. Identification on \(\psi_j\) means that
\[
\neg\exists\theta^0\exists\theta^1
\left[
\mathsf C_{\mathcal G}(\theta^0;\mathbf p_{\mathrm{obs}})
\wedge\mathsf C_{\mathcal G}(\theta^1;\mathbf p_{\mathrm{obs}})
\wedge q_{a,y}(\theta^0)\ne q_{a,y}(\theta^1)
\right]
\tag{2}
\]
holds. Nonemptiness follows from (1), so the relation
\[
\mathsf V_{a,y}(\mathbf p_{\mathrm{obs}},z)
\Longleftrightarrow
\exists\theta\left[
\mathsf C_{\mathcal G}(\theta;\mathbf p_{\mathrm{obs}})
\wedge z=q_{a,y}(\theta)
\right]
\tag{3}
\]
has exactly one solution. Substitution of \(\theta_{\mathcal M}\) shows that this solution is
\[
z_{a,y}=Q_a^{\mathcal M}(y).
\tag{4}
\]

If the cell retains the value graph (3), its functional is (4) by definition. If it carries a rational annotation \(N/D\), \(\text{def:tcf-procedure}\) accepts that annotation only after verifying
\[
\forall\mathbf p\forall z\,
\left[
\psi_j(\mathbf p)\wedge\mathsf V_{a,y}(\mathbf p,z)
\Longrightarrow D(\mathbf p)\ne0\wedge D(\mathbf p)z=N(\mathbf p)
\right].
\tag{5}
\]
At \((\mathbf p_{\mathrm{obs}},z_{a,y})\), (5) licenses division and yields \(N(\mathbf p_{\mathrm{obs}})/D(\mathbf p_{\mathrm{obs}})=z_{a,y}\). Thus in either representation
\[
f_{a,j}(P_{\mathrm{obs}})(y)=Q_a^{\mathcal M}(y).
\tag{6}
\]
This argument uses the whole-cell implication, not Markovness of a discovered raw trace.

The choices of \(\mathcal M,a,y\) were arbitrary. Since \(\mathcal X_Y\) is finite, substitute (6) term by term in \(\text{def:causal-query}\):
\[
\begin{aligned}
&\sum_y s_Y(y)
\{f_{a_1,j}(P_{\mathrm{obs}})(y)-f_{a_0,j}(P_{\mathrm{obs}})(y)\}\\
&\qquad=\sum_y s_Y(y)
\{Q_{a_1}^{\mathcal M}(y)-Q_{a_0}^{\mathcal M}(y)\}
=\tau_{s_Y}^{\mathcal M}.
\end{aligned}
\tag{7}
\]
The proof treats the supplied scores as arbitrary real coefficients and therefore requires no finite encoding of the population point.

%%% FIELD tcf-termination-proof
For a raw trace state \(\kappa\), let \(r(\kappa)\) be the number of random vertices and let \(d(c)\) be the number of assigned missingness indicators. Put
\[
\mu(\kappa)=r(\kappa)+|\mathcal R|-d(c).
\tag{1}
\]
This nonnegative integer is initially at most \(2n\). Each strict ancestral marginalization, proper district-child restriction, or fixing strictly decreases the random set, while each delayed selection strictly increases \(d(c)\). Hence every enumerated operation word has length at most \(2n\), independently of any confluence or branch merge.

There are at most \(2^n\) initial observed margins and
\[
c=\prod_{R_Z\in\mathcal R}(|\mathcal X_{R_Z}|+1)
\tag{2}
\]
partial contexts. The two subset operations contribute at most \(2^{n+1}\) symbols, fixing contributes at most \(n(\ell+1)\), selection contributes at most \(c\), and the terminal marker contributes one. Therefore
\[
L_{\mathrm{op}}=2^{n+1}+n(\ell+1)+c+1,
\qquad
B_{\mathrm{state}}=2^n c\sum_{j=0}^{2n}L_{\mathrm{op}}^j.
\tag{3}
\]
A fixed initial trace and fixed word determine every component of the literal tuple \((\mathcal H,F,c,q,\pi,\Lambda)\). Since the corrected procedure never merges distinct words, (3) directly bounds all visited raw states.

Write \(h=|\mathcal E^{\leftrightarrow}(\mathcal G)|+|\mathcal V|+|\mathcal R|\) and \(g=|\mathcal T_{\mathcal G}|\). A doubled coordinate test has at most
\[
p_{\mathrm{coord}}=2\{g+hK_{\mathrm{ub}}(\mathcal G)\}+|\mathcal X_{\mathcal O}|
\tag{4}
\]
quantified-plus-free variables, and the direct target test has at most
\[
p_{\tau}=p_{\mathrm{coord}}+|\mathcal X_Y|.
\tag{5}
\]
The encoding length of either base input is at most \(s\). Real-closed-field quantifier elimination and finite sign-invariant CAD, recorded in \(\text{lem:real-closed-field-quantifier-elimination}\) and \(\text{lem:finite-sign-invariant-cad}\), therefore terminate, and the definition of \(\operatorname{CAD}\) bounds the two base calls by
\[
C_{\mathrm{base}}
=\operatorname{CAD}(p_{\mathrm{coord}},s)
+\operatorname{CAD}(p_{\tau},s).
\tag{6}
\]

Each raw state supplies at most one candidate per coordinate and coordinate cell and at most one candidate per target cell. Since the number of cells is included in the corresponding \(\operatorname{CAD}\) bound, the number of later verification sentences is at most
\[
N_{\mathrm{check}}
=B_{\mathrm{state}}
\left\{|\mathcal X_A||\mathcal X_Y|
\operatorname{CAD}(p_{\mathrm{coord}},s)
+\operatorname{CAD}(p_{\tau},s)\right\}.
\tag{7}
\]
After denominators are cleared, a verification sentence uses at most one response vector, the value variable \(z\), the observed coordinates, and, for a target sentence, the score coordinates. It therefore has at most \(p_{\max}+1\) variables and length at most \(s\), so each such sentence costs at most \(\operatorname{CAD}(p_{\max}+1,s)\). Summing (3), (6), and (7) gives
\[
B_{\mathrm{tot}}
=B_{\mathrm{state}}+C_{\mathrm{base}}
+N_{\mathrm{check}}\operatorname{CAD}(p_{\max}+1,s).
\tag{8}
\]

Both partitions contain finitely many integer-polynomial sign tests. The interface \(\operatorname{Exact}(P_{\mathrm{obs}})\) decides the coordinate signs, and \(\operatorname{Exact}(P_{\mathrm{obs}},s_Y)\) decides the joint observed-and-score signs. The fixed CAD procedure then performs finitely many root-isolation and sign-evaluation steps and supplies the stated Thom representations over the relevant generated ordered field. These specialization steps are finite but are defined only for the stated exact interfaces. Equations (1)--(8) prove termination and the explicit bound, with no polynomial-time assertion.

%%% FIELD qe-compiler-proof
Because \(\omega=(a,y,\theta^0,\theta^1)\) satisfies \(\mathsf W_j\),
\[
\mathsf C_{\mathcal G}(\theta^i;P_{\mathrm{obs}})\quad(i=0,1),
\qquad
q_{a,y}(\theta^0)\ne q_{a,y}(\theta^1).
\tag{1}
\]
The established \(\text{lem:negative-cell-model-pair}\) decodes these vectors into models satisfying
\[
\begin{gathered}
\mathcal M_i\in\mathfrak F(\mathcal G,P_{\mathrm{obs}}),
\qquad
P_{\mathcal M_i}(\mathcal O)=P_{\mathrm{obs}},\\
Q_a^{\mathcal M_i}(y)=q_{a,y}(\theta^i),
\qquad i=0,1.
\end{gathered}
\tag{2}
\]
That lemma also gives strict positivity and the product over the individual bidirected-edge and private factors. Equations (1)--(2) imply
\[
Q_a^{\mathcal M_0}(y)\ne Q_a^{\mathcal M_1}(y).
\tag{3}
\]

The retained history \(\mathcal T\) and the graph are finite. Hence the fixed prefix and subset order in \(\text{def:certificate-compiler}\) contains finitely many triples \((\kappa_*,D,D')\). Ancestor, district, root, inclusion, and treatment-incidence tests are finite graph computations. Every displayed trace equality is rational, and its recorded support makes each denominator nonzero; multiplying by their finite product is therefore equivalence-preserving and produces a polynomial sign test.

By \(\text{def:exact-observed-specialization}\), the coordinates of \(\theta^0,\theta^1\) have Thom encodings over the real closure of the ordered field generated by \((P_{\mathrm{obs}}(o))_o\). Signed-subresultant evaluation reduces each candidate polynomial comparison to finitely many signs at those coordinates, and \(\operatorname{Exact}(P_{\mathrm{obs}})\) decides them. Thus every candidate test terminates. A finite deterministic enumeration therefore returns the first triple passing every conjunct of \(\Delta\), or reaches its end and returns \(\operatorname{exhausted}\). Equations (2)--(3) hold independently of that tag. Since no step proves \(\operatorname{Inv}(\kappa_*)\), a forest realization, or membership in \(\operatorname{CH}(\mathcal G,A,Y)\), the final non-assertion is also exact.

%%% FIELD pointwise-statement
For a finite positive regular-maximal \(\mathcal G\) and supplied \(A,Y\), the trace-independent base phase of \(\operatorname{TCF}\) constructs the finite coordinate partition \(\mathscr P_{\mathcal G,A,Y}\) of \(\{\mathbf p:\exists\theta\,\mathsf C_{\mathcal G}(\theta;\mathbf p)\}\) and the finite direct-target partition \(\mathscr P^{\tau}_{\mathcal G,A,Y}\) of that feasible projection times \(\mathbb R^{\mathcal X_Y}\). Raw trace enumeration only proposes annotations, and an annotation is retained only after whole-cell verification. For every compatible finite/full-positive \(P_{\mathrm{obs}}\), every \((a,y)\), and every supplied real score map \(s_Y\), its coordinate cell carries \(f_{a,j}(P_{\mathrm{obs}})(y)\) if and only if \(\mathcal M\mapsto Q_a^{\mathcal M}(y)\) is constant on \(\mathfrak F(\mathcal G,P_{\mathrm{obs}})\), and then
\[
f_{a,j}(P_{\mathrm{obs}})(y)=Q_a^{\mathcal M}(y)
\quad\text{for every }\mathcal M\in\mathfrak F(\mathcal G,P_{\mathrm{obs}}).
\]
Its target cell at \((P_{\mathrm{obs}},s_Y)\) carries \(f_{\tau,k}(P_{\mathrm{obs}};s_Y)\) if and only if \(\mathcal M\mapsto\tau_{s_Y}^{\mathcal M}\) is constant on the same fiber, and then
\[
f_{\tau,k}(P_{\mathrm{obs}};s_Y)=\tau_{s_Y}^{\mathcal M}
\quad\text{for every }\mathcal M\in\mathfrak F(\mathcal G,P_{\mathrm{obs}}).
\]
This direct equivalence remains valid when some interventional coordinates vary but their score contrast is constant. If a coordinate or target is not constant, its cell carries \(\mathsf W_j\) or \(\mathsf W_{\tau}\), and without any exact encoding there exist strictly positive individual-edge-factorized models \(\mathcal M_0,\mathcal M_1\in\mathfrak F(\mathcal G,P_{\mathrm{obs}})\subseteq\mathfrak M_{\mathrm{lm}}(\mathcal G)\) with the same observed proxy law and different values of the corresponding coordinate or score contrast. Under \(\operatorname{Exact}(P_{\mathrm{obs}})\), coordinate witnesses are represented and \(\operatorname{Compile}_{\mathrm{col}}\) terminates with a checked divergence record or an exhaustion tag; under \(\operatorname{Exact}(P_{\mathrm{obs}},s_Y)\), target witnesses are represented. If every coordinate is constant at \(a_0\) and \(a_1\), then
\[
\tau_{s_Y}^{\mathcal M}
=\sum_{y\in\mathcal X_Y}s_Y(y)
\{f_{a_1,j}(P_{\mathrm{obs}})(y)-f_{a_0,j}(P_{\mathrm{obs}})(y)\}.
\]
This is pointwise semialgebraic completeness on the finite/full-positive subclass. It is not a graphwise colored-recursion or colored-hedge completeness theorem.

%%% FIELD pointwise-proof
By \(\text{thm:tcf-termination}\), both trace-independent base constructions and the subsequent finite raw-trace annotation search terminate. By \(\text{lem:semialgebraic-fiber-dichotomy}\), the coordinate cells form a finite partition of the feasible observed-law projection and the target cells form a finite partition of its product with the real score space. The construction in \(\text{def:tcf-procedure}\) fixes both partitions before raw traces are enumerated and accepts an annotation only after its displayed whole-cell implication has been decided true. Hence trace discovery changes neither partition.

Fix compatible \(P_{\mathrm{obs}}\), a real score map \(s_Y\), and write
\[
\mathbf p_{\mathrm{obs}}=(P_{\mathrm{obs}}(o))_o,
\qquad
\mathbf s_Y=(s_Y(y))_y.
\tag{1}
\]
The dichotomy lemma proves
\[
\mathsf I_{a,y}(\mathcal G;\mathbf p_{\mathrm{obs}})
\Longleftrightarrow
\left[\mathcal M\mapsto Q_a^{\mathcal M}(y)
\text{ is constant on }\mathfrak F(\mathcal G,P_{\mathrm{obs}})\right]
\tag{2}
\]
and
\[
\mathsf I_{\tau}(\mathcal G;\mathbf p_{\mathrm{obs}},\mathbf s_Y)
\Longleftrightarrow
\left[\mathcal M\mapsto\tau_{s_Y}^{\mathcal M}
\text{ is constant on }\mathfrak F(\mathcal G,P_{\mathrm{obs}})\right].
\tag{3}
\]

When (2) holds, \(\text{thm:tcf-soundness}\) identifies the unique value graph, and it also identifies every accepted rational annotation because its whole-cell verification supplies the required nonzero denominator. Therefore
\[
f_{a,j}(P_{\mathrm{obs}})(y)=Q_a^{\mathcal M}(y)
\quad\text{for all }\mathcal M\in\mathfrak F(\mathcal G,P_{\mathrm{obs}}).
\tag{4}
\]
Conversely, a common value in (4) makes the coordinate constant, and (2) holds. This proves the coordinate equivalence.

For the direct target, \(\text{thm:response-fiber-equivalence}\) gives, for every feasible \(\theta_{\mathcal M}\),
\[
t_{\mathbf s_Y}(\theta_{\mathcal M})
=\sum_y s_Y(y)\{Q_{a_1}^{\mathcal M}(y)-Q_{a_0}^{\mathcal M}(y)\}
=\tau_{s_Y}^{\mathcal M}.
\tag{5}
\]
When (3) holds, the target value relation is therefore nonempty and single-valued at \(\tau_{s_Y}^{\mathcal M}\). If a rational target annotation \(N_{\tau}/D_{\tau}\) is retained, its whole-cell test gives
\[
D_{\tau}(\mathbf p_{\mathrm{obs}},\mathbf s_Y)\ne0,
\qquad
D_{\tau}(\mathbf p_{\mathrm{obs}},\mathbf s_Y)\tau_{s_Y}^{\mathcal M}
=N_{\tau}(\mathbf p_{\mathrm{obs}},\mathbf s_Y),
\tag{6}
\]
so division is licensed and the annotation has the same value. Thus
\[
f_{\tau,k}(P_{\mathrm{obs}};s_Y)=\tau_{s_Y}^{\mathcal M}
\quad\text{for all }\mathcal M\in\mathfrak F(\mathcal G,P_{\mathrm{obs}}),
\tag{7}
\]
exactly when (3) holds. Equation (5) proves this without coordinatewise constancy.

If a coordinate test is false, the coordinate clause of the dichotomy lemma supplies \(\mathsf W_j\) and two strictly positive individual-edge-factorized models in the same observed fiber with different coordinate values. If the target test is false, its target clause supplies \(\mathsf W_{\tau}\) and models satisfying
\[
P_{\mathcal M_0}(\mathcal O)=P_{\mathcal M_1}(\mathcal O)=P_{\mathrm{obs}},
\qquad
\tau_{s_Y}^{\mathcal M_0}\ne\tau_{s_Y}^{\mathcal M_1}.
\tag{8}
\]
No exact encoding is needed for these existential statements. By \(\text{prop:ambient-nonrecovery-transfer}\), the same pairs prove nonidentification in the larger published ambient fiber. Under the appropriate exact interface, \(\text{def:exact-observed-specialization}\) represents the pair; for a coordinate pair, \(\text{thm:qe-failure-compiler}\) terminates with its diagnostic tag. That tag is not used as a hedge certificate.

Finally, if every coordinate at \(a_0\) and \(a_1\) is constant, finiteness of \(\mathcal X_Y\), equation (4), and \(\text{def:causal-query}\) give
\[
\begin{aligned}
\tau_{s_Y}^{\mathcal M}
&=\sum_y s_Y(y)\{Q_{a_1}^{\mathcal M}(y)-Q_{a_0}^{\mathcal M}(y)\}\\
&=\sum_y s_Y(y)
\{f_{a_1,j}(P_{\mathrm{obs}})(y)-f_{a_0,j}(P_{\mathrm{obs}})(y)\}.
\end{aligned}
\tag{9}
\]
Equations (1)--(9) prove every pointwise assertion. They use neither a graph-only recursion-failure theorem nor colored-hedge realization.

%%% FIELD two-comparator-current
On the common finite, strictly positive regular-maximal labeled-missingness subclass, TCF is a strict class-relative extension of the de Aguas Proposition 4 criterion and an embedding of the Mokhtarian observed-root theorem. More precisely, for every finite positive regular-maximal \(\mathcal G\), compatible \(P_{\mathrm{obs}}\), treatment value \(a\), outcome value \(y\), and \(\xi\in\operatorname{DA4}(\mathcal G,A,Y)\), the unique TCF cell \(\psi_j\in\mathscr P_{\mathcal G,A,Y}\) carries \[f_{a,j}(P_{\mathrm{obs}})(y)=f^{\mathrm{DA4}}_{a,\xi}(P_{\mathrm{obs}})(y).\] Whenever \(\operatorname{RC}(\mathcal G,A,Y)\) holds, the same cell carries \(f_{a,j}(P_{\mathrm{obs}})(y)=f_a^{\mathrm{RC}}(P_{\mathrm{obs}})(y)\). The de Aguas containment is strict on \(\mathfrak G_{\mathrm{EIPFD}}\): for every \(\mathcal G(H,m)\) and every compatible \(P_{\mathrm{obs}}\), \(\operatorname{DA4}(\mathcal G(H,m),A,Y)=\varnothing\), while the TCF cell carries \[\sum_hP_{\mathrm{obs}}(h)\sum_bP_{\mathrm{obs}}(b\mid a,h,R_W=0)\sum_{a'}P_{\mathrm{obs}}(y\mid b,a',h,R_W=0)P_{\mathrm{obs}}(a'\mid h,R_W=0),\] where \(b\) ranges over \(B_1\). The arrow \(C\to R_W\) shows that the observed-root premise is inapplicable to this family; it is not a strictness claim about the complete published root-control domain. For \(H=\{C\}\) and \(m=1\), the displayed functional is exactly \(f_a^\star(P_{\mathrm{obs}})(y)\) of \(\text{prop:eipfd-star-worked}\).

%%% FIELD two-comparator-proposed
On the common finite, strictly positive regular-maximal labeled-missingness subclass, TCF is a strict class-relative extension of the de Aguas Proposition 4 criterion and an embedding of the Mokhtarian observed-root theorem. More precisely, for every finite positive regular-maximal \(\mathcal G\), compatible \(P_{\mathrm{obs}}\), treatment value \(a\), outcome value \(y\), and \(\xi\in\operatorname{DA4}(\mathcal G,A,Y)\), the unique TCF cell \(\psi_j\in\mathscr P_{\mathcal G,A,Y}\) carries \[f_{a,j}(P_{\mathrm{obs}})(y)=f^{\mathrm{DA4}}_{a,\xi}(P_{\mathrm{obs}})(y).\] Whenever \(\operatorname{RC}(\mathcal G,A,Y)\) holds, the same cell carries \(f_{a,j}(P_{\mathrm{obs}})(y)=f_a^{\mathrm{RC}}(P_{\mathrm{obs}})(y)\). The de Aguas containment is strict on \(\mathfrak G_{\mathrm{EIPFD}}\): for every \(\mathcal G(H,m)\) and every compatible \(P_{\mathrm{obs}}\), \(\operatorname{DA4}(\mathcal G(H,m),A,Y)=\varnothing\), while the TCF cell carries \[\sum_hP_{\mathrm{obs}}(h)\sum_bP_{\mathrm{obs}}(b\mid a,h,R_W=0)\sum_{a'}P_{\mathrm{obs}}(y\mid b,a',h,R_W=0)P_{\mathrm{obs}}(a'\mid h,R_W=0),\] where \(b\) ranges over \(B_1\). The arrow \(C\to R_W\) shows that the observed-root premise is inapplicable to this family; it is not a strictness claim about the complete published root-control domain. For \(H=\{C\}\) and \(m=1\), the displayed functional is exactly \(f_a^\star(P_{\mathrm{obs}})(y)\) of \(\text{prop:eipfd-star-worked}\). Relative to the general finite discrete response-function polynomial-programming framework of Duarte et al., the response-polynomial/CAD layer is only an lm-specific constrained specialization: \(\mathsf C_{\mathcal G}\) restricts the feasible set to one label-compatible table tuple and mutually independent simplices for the individual bidirected-edge and private factors. No strict-extension, non-subsumption, or sharp-bound claim relative to Duarte et al. is made.

%%% FIELD response-fiber-statement
For every finite positive regular-maximal \(\mathcal G\) and every observed law \(P_{\mathrm{obs}}\), the polynomial presentation is exactly the compatible lm-SCM fiber in both directions. First, for every \(\mathcal M\in\mathfrak F(\mathcal G,P_{\mathrm{obs}})\), finite support reduction of each mutually independent individual-edge and private latent factor produces \(\theta_{\mathcal M}\) such that \(\mathsf C_{\mathcal G}(\theta_{\mathcal M};P_{\mathrm{obs}})\), \(p_{v,r}(\theta_{\mathcal M})=P_{\mathcal M}(v,r)\), and \(q_{a,y}(\theta_{\mathcal M})=Q_a^{\mathcal M}(y)\) for every \((v,r,a,y)\). Conversely, every \(\theta\) satisfying \(\mathsf C_{\mathcal G}(\theta;P_{\mathrm{obs}})\) decodes, using the unique selected table and the product law \(\prod_h\rho_h\), to a model \(\mathcal M_{\theta}\in\mathfrak F(\mathcal G,P_{\mathrm{obs}})\) with the same full-data cells and \(Q_a^{\mathcal M_{\theta}}(y)=q_{a,y}(\theta)\).

%%% FIELD tcf-soundness-statement
For every identifying cell \(\psi_j(\mathbf p)\in\mathscr P_{\mathcal G,A,Y}\), every compatible law \(P_{\mathrm{obs}}\) satisfying \(\psi_j((P_{\mathrm{obs}}(o))_o)\), every \(\mathcal M\in\mathfrak F(\mathcal G,P_{\mathrm{obs}})\), every \(a\in\mathcal X_A\), and every \(y\in\mathcal X_Y\), the cell's functional satisfies \[f_{a,j}(P_{\mathrm{obs}})(y)=Q_a^{\mathcal M}(y).\] Consequently, \(\sum_y s_Y(y)\{f_{a_1,j}(P_{\mathrm{obs}})(y)-f_{a_0,j}(P_{\mathrm{obs}})(y)\}=\tau_{s_Y}\) throughout that fiber. This is a mathematical specialization statement and does not require a finite encoding of the population coordinates.

%%% FIELD tcf-termination-statement
Let the quantities in \(\text{def:termination-budget}\) be fixed. For every finite \(\mathcal G\) and finite alphabets, the expression-discovery phase of \(\operatorname{TCF}(\mathcal G,A,Y,P_{\mathrm{obs}})\) visits at most
\[
B_{\mathrm{state}}=2^n c\sum_{j=0}^{2n}\bigl\{2^{n+1}+n(\ell+1)+c+1\bigr\}^j
\]
literal complete trace states \(\kappa=(\mathcal H,F,c,q,\pi,\Lambda)\), counting every fixing-selection interleaving and provenance-bearing history. Independently of those traces, the coordinate and direct-target base CADs construct \(\mathscr P_{\mathcal G,A,Y}\) and \(\mathscr P^{\tau}_{\mathcal G,A,Y}\) within
\[
C_{\mathrm{base}}=\operatorname{CAD}(p_{\mathrm{coord}},s)+\operatorname{CAD}(p_{\tau},s)
\]
projection, root-isolation, sign-evaluation steps, and output cells. At most
\[
N_{\mathrm{check}}=B_{\mathrm{state}}\left\{|\mathcal X_A||\mathcal X_Y|\operatorname{CAD}(p_{\mathrm{coord}},s)+\operatorname{CAD}(p_{\tau},s)\right\}
\]
coordinate and target whole-cell candidate-equality sentences are subsequently decided, and the total work is bounded by
\[
B_{\mathrm{tot}}=B_{\mathrm{state}}+C_{\mathrm{base}}+N_{\mathrm{check}}\operatorname{CAD}(p_{\max}+1,s).
\]
Given \(\operatorname{Exact}(P_{\mathrm{obs}})\), coordinate-cell selection and represented value or witness extraction terminate after finitely many exact sign decisions. Given \(\operatorname{Exact}(P_{\mathrm{obs}},s_Y)\), the same is true for the direct-target cell over the ordered field generated by the observed and score coordinates. No point computation is claimed for arbitrary unencoded reals, and no polynomial-time claim is made.

%%% FIELD two-comparator-proof
Let \(P_{\mathrm{obs}}\) be compatible with the common finite positive class, fix \((a,y)\), and let \(\psi_j\) be its unique TCF coordinate cell. For \(\xi\in\operatorname{DA4}(\mathcal G,A,Y)\), \(\text{lem:da4-recovery}\) gives
\[
f^{\mathrm{DA4}}_{a,\xi}(P_{\mathrm{obs}})(y)
=Q_a^{\mathcal M}(y)
\quad\text{for every }\mathcal M\in\mathfrak F(\mathcal G,P_{\mathrm{obs}}).
\tag{1}
\]
The left side depends only on the observed law, so the coordinate is constant on the fiber. The semialgebraic dichotomy and \(\text{thm:tcf-soundness}\) therefore imply
\[
f_{a,j}(P_{\mathrm{obs}})(y)
=Q_a^{\mathcal M}(y)
=f^{\mathrm{DA4}}_{a,\xi}(P_{\mathrm{obs}})(y).
\tag{2}
\]

If \(\operatorname{RC}(\mathcal G,A,Y)\) holds, \(\text{lem:root-control-context-recovery}\) analogously gives
\[
f_a^{\mathrm{RC}}(P_{\mathrm{obs}})(y)
=Q_a^{\mathcal M}(y)
\quad\text{throughout the fiber},
\tag{3}
\]
and the same constancy and soundness argument yields
\[
f_{a,j}(P_{\mathrm{obs}})(y)=f_a^{\mathrm{RC}}(P_{\mathrm{obs}})(y).
\tag{4}
\]

For strictness of the de Aguas containment, fix \(\mathcal G(H,m)\in\mathfrak G_{\mathrm{EIPFD}}\). The bidirected edge \(A\leftrightarrow Y\) survives deletion of directed arrows out of \(A\), and a length-one path has no nonendpoint that a proposed adjustment set can block. Hence condition \(\mathrm{(v)}\) fails for every proposed de Aguas certificate, so
\[
\operatorname{DA4}(\mathcal G(H,m),A,Y)=\varnothing.
\tag{5}
\]
On the other hand, \(\text{lem:eipfd-frontdoor-identification}\) gives, for every model in a compatible fiber,
\[
\begin{aligned}
Q_a^{\mathcal M}(y)
&=\sum_hP_{\mathrm{obs}}(h)
\sum_{b\in\mathcal X_{B_1}}P_{\mathrm{obs}}(b\mid a,h,R_W=0)\\
&\qquad\times\sum_{a'\in\mathcal X_A}
P_{\mathrm{obs}}(y\mid b,a',h,R_W=0)
P_{\mathrm{obs}}(a'\mid h,R_W=0).
\end{aligned}
\tag{6}
\]
Strict positivity supplies every conditional denominator in (6). Its right side is an observed-law functional, so the coordinate is constant. The dichotomy and soundness prove that the TCF cell carries (6). Equations (5)--(6) prove strict class-relative containment. Every family member contains \(C\to R_W\), so the root-control premise is inapplicable; this does not compare performance on the rest of the root-control domain. For \(H=\{C\}\) and \(m=1\), \(\text{prop:eipfd-star-worked}\) identifies (6) with \(f_a^\star(P_{\mathrm{obs}})(y)\).

Finally, \(\text{lem:duarte-response-polynomial-scope}\) supplies the published scope of finite response-function polynomial programming. Unfolding \(\text{def:response-polynomial-fiber}\), the TCF feasibility set uses response tables but adds the equations selecting one \(g_*\in\mathcal T_{\mathcal G}\), the label-imposed table identifications, and one independent simplex for each \(h\in\mathcal E^{\leftrightarrow}(\mathcal G)\sqcup\mathcal V\sqcup\mathcal R\). Thus the present response-CAD layer is obtained by imposing lm-specific constraints inside that response-polynomial paradigm. This proves the specialization statement and, because no reverse containment or performance comparison was used, proves the three explicit nonclaims.
